
\Needspace{15\baselineskip}
\subsubsection{Evaluaciones internas de firmas explícitas y contraste posterior}

La tabla siguiente publica de forma legible las diecisiete coordenadas HMT--MD ya
calculadas ---once elementales y seis compuestas--- y las confronta con la
edición 2026 del conjunto de datos del Particle Data Group, con fecha de corte
editorial de 15 de enero de 2026 \cite{PDG2026}. No constituyen otra ontología
distinta de la capa operatoria anterior. El símbolo \(\dagger\)
identifica una evaluación obtenida al aplicar el mapa de construcción a una firma
explícita. En las once filas elementales la firma procede de la ruta nominal
determinada sin masas; en las seis compuestas procede del registro de ligadura
declarado. Cada cifra depende de esa firma, del lector bidireccional y de la carta
dimensional común; el conjunto de datos externo interviene sólo en la comparación
y no selecciona sabor, generación, color ni ruta.


Para leptones cargados y bosones masivos se muestra la diferencia relativa entre
las dos coordenadas centrales y, cuando el conjunto de datos aporta una incertidumbre
bilateral, su cociente por esa incertidumbre. En los quarks se publica la
diferencia aritmética entre valores centrales, pero se reserva el término
\emph{residual metrológico}: éste exige que la masa interna y la referencia externa
usen la misma definición, el mismo esquema de renormalización y la misma escala.
Las masas ligeras se leen en la convención \(\overline{\mathrm{MS}}\) a
\(2\,\mathrm{GeV}\); las de \(c\) y \(b\), en
\(\mu=\overline m_q\); la masa superior procede de reconstrucción cinemática.

\begingroup
\setlength{\tabcolsep}{4.5pt}
\renewcommand{\arraystretch}{1.28}
\begin{longtable}{@{}>{\raggedright\arraybackslash}p{.12\textwidth}>{\raggedright\arraybackslash}p{.27\textwidth}>{\raggedright\arraybackslash}p{.27\textwidth}>{\raggedright\arraybackslash}p{.28\textwidth}@{}}
\caption{Evaluaciones internas de firmas explícitas y contraste metrológico posterior.}\label{tab:masas-evaluaciones-firmas-explicitas}\\
\toprule
\textbf{Estado} & \textbf{Coordenada de la firma} & \textbf{Referencia física externa} & \textbf{Alcance} \\
\midrule
\endfirsthead
\toprule
\textbf{Estado} & \textbf{Coordenada de la firma} & \textbf{Referencia física externa} & \textbf{Alcance} \\
\midrule
\endhead
\midrule
\multicolumn{4}{r}{\footnotesize Continúa en la página siguiente.}\\
\endfoot
\bottomrule
\endlastfoot
% HMT-MASS-EXPLICIT-SIGNATURE:SM-LC-2
$\mu^-/\mu^+$ & $105.658360\allowbreak{}294435\allowbreak{}829303\allowbreak{}4\;\mathrm{MeV}/c^2$\textsuperscript{\(\dagger\)}; torre E3, firma (42,-3,8,0,2) & $105.6583755\pm0.0000023\;\mathrm{MeV}/c^2$; masa leptónica de reposo; promedio PDG y conversión metrológica declarada & Diferencia relativa: $-0.143912\allowbreak{}530\;\mathrm{ppm}$\textsuperscript{\(\dagger\)}; $-6.611114\,\sigma$. La cifra HMT es la evaluación exacta de la firma declarada. \\
% HMT-MASS-EXPLICIT-SIGNATURE:SM-LC-3
$\tau^-/\tau^+$ & $1776.860917\allowbreak{}631432\allowbreak{}295595\;\mathrm{MeV}/c^2$\textsuperscript{\(\dagger\)}; torre E3, firma (64,0,25,-1,6) & $1776.93\pm0.09\;\mathrm{MeV}/c^2$; promedio directo PDG & Diferencia relativa: $-38.877371\allowbreak{}966\;\mathrm{ppm}$\textsuperscript{\(\dagger\)}; $-0.813547\,\sigma$. La cifra HMT es la evaluación exacta de la firma declarada. \\
% HMT-MASS-EXPLICIT-SIGNATURE:SM-Q-U1
$u/\bar u$ & $2.165924\allowbreak{}536173\allowbreak{}907\;\mathrm{MeV}/c^2$\textsuperscript{\(\dagger\)}; carta angular, firma nominal de ruta (11,7) & $2.16\pm0.07\;\mathrm{MeV}/c^2$; masa $\overline{\mathrm{MS}}$ a $\mu=2\,\mathrm{GeV}$; 90\% C.L. & Diferencia entre valores centrales: $2742.840821\allowbreak{}253\;\mathrm{ppm}$\textsuperscript{\(\dagger\)}. No es un residual metrológico hasta fijar el mismo esquema y la misma escala. \\
% HMT-MASS-EXPLICIT-SIGNATURE:SM-Q-D1
$d/\bar d$ & $4.679550\allowbreak{}162948\allowbreak{}874\;\mathrm{MeV}/c^2$\textsuperscript{\(\dagger\)}; carta angular, firma nominal de ruta (17,8) & $4.70\pm0.07\;\mathrm{MeV}/c^2$; masa $\overline{\mathrm{MS}}$ a $\mu=2\,\mathrm{GeV}$; 90\% C.L. & Diferencia entre valores centrales: $-4351.029159\allowbreak{}814\;\mathrm{ppm}$\textsuperscript{\(\dagger\)}. No es un residual metrológico hasta fijar el mismo esquema y la misma escala. \\
% HMT-MASS-EXPLICIT-SIGNATURE:SM-Q-U2
$c/\bar c$ & $1.270500\allowbreak{}838641\allowbreak{}911\;\mathrm{GeV}/c^2$\textsuperscript{\(\dagger\)}; carta angular, firma nominal de ruta (61,8) & $1.2729\pm0.0045\;\mathrm{GeV}/c^2$; masa $\overline{\mathrm{MS}}$ evaluada en $\mu=\overline{m}_q$; 90\% C.L. & Diferencia entre valores centrales: $-1884.799558\allowbreak{}558\;\mathrm{ppm}$\textsuperscript{\(\dagger\)}. No es un residual metrológico hasta fijar el mismo esquema y la misma escala. \\
% HMT-MASS-EXPLICIT-SIGNATURE:SM-Q-D2
$s/\bar s$ & $92.940093\allowbreak{}907845\allowbreak{}64\;\mathrm{MeV}/c^2$\textsuperscript{\(\dagger\)}; carta angular, firma nominal de ruta (41,-3) & $92.9\pm0.7\;\mathrm{MeV}/c^2$; masa $\overline{\mathrm{MS}}$ a $\mu=2\,\mathrm{GeV}$; 90\% C.L. & Diferencia entre valores centrales: $431.581354\allowbreak{}635\;\mathrm{ppm}$\textsuperscript{\(\dagger\)}. No es un residual metrológico hasta fijar el mismo esquema y la misma escala. \\
% HMT-MASS-EXPLICIT-SIGNATURE:SM-Q-U3
$t/\bar t$ & $172.597479\allowbreak{}544019\allowbreak{}1\;\mathrm{GeV}/c^2$\textsuperscript{\(\dagger\)}; carta angular, firma nominal de ruta (100,-1) & $172.60\pm0.27\;\mathrm{GeV}/c^2$; masa reconstruida directamente a partir de la cinemática de sucesos & Diferencia entre valores centrales: $-14.602873\allowbreak{}585\;\mathrm{ppm}$\textsuperscript{\(\dagger\)}. No es un residual metrológico hasta fijar el mismo esquema y la misma escala. \\
% HMT-MASS-EXPLICIT-SIGNATURE:SM-Q-D3
$b/\bar b$ & $4.190119\allowbreak{}763299\allowbreak{}790\;\mathrm{GeV}/c^2$\textsuperscript{\(\dagger\)}; carta angular, firma nominal de ruta (71,-5) & $4.186\pm0.006\;\mathrm{GeV}/c^2$; masa $\overline{\mathrm{MS}}$ evaluada en $\mu=\overline{m}_q$; 90\% C.L. & Diferencia entre valores centrales: $984.176612\allowbreak{}467\;\mathrm{ppm}$\textsuperscript{\(\dagger\)}. No es un residual metrológico hasta fijar el mismo esquema y la misma escala. \\
% HMT-MASS-EXPLICIT-SIGNATURE:SM-GA-W
$W^\pm$ & $80.378964\allowbreak{}909704\allowbreak{}547024\allowbreak{}86\;\mathrm{GeV}/c^2$\textsuperscript{\(\dagger\)}; torre E3, firma (94,-1,0,-2,4) & $80.3625\pm0.0077\;\mathrm{GeV}/c^2$; promedio de masa de la instantánea PDG sellada & Diferencia relativa: $204.882995\allowbreak{}234\;\mathrm{ppm}$\textsuperscript{\(\dagger\)}; $2.138299\,\sigma$. La cifra HMT es la evaluación exacta de la firma declarada. \\
% HMT-MASS-EXPLICIT-SIGNATURE:SM-GA-Z
$Z^0$ & $91.187533\allowbreak{}444313\allowbreak{}407844\allowbreak{}85\;\mathrm{GeV}/c^2$\textsuperscript{\(\dagger\)}; torre E3, firma (95,-1,-11,0,-4) & $91.1879\pm0.0020\;\mathrm{GeV}/c^2$; promedio de masa de la instantánea PDG sellada & Diferencia relativa: $-4.019784\allowbreak{}276\;\mathrm{ppm}$\textsuperscript{\(\dagger\)}; $-0.186362\,\sigma$. La cifra HMT es la evaluación exacta de la firma declarada. \\
% HMT-MASS-EXPLICIT-SIGNATURE:SM-H-1
Higgs $H$ & $125.292466\allowbreak{}042536\allowbreak{}133\;\mathrm{GeV}/c^2$\textsuperscript{\(\dagger\)}; carta angular, firma nominal de ruta (97,9) & $125.13\pm0.11\;\mathrm{GeV}/c^2$; promedio de masa de la instantánea PDG sellada & Diferencia relativa: $1298.378027\allowbreak{}140\;\mathrm{ppm}$\textsuperscript{\(\dagger\)}; $1.454206\,\sigma$. La cifra HMT es la evaluación exacta de la firma declarada. \\
% HMT-MASS-EXPLICIT-SIGNATURE:E3B-002
$\pi^0$\par firma $(44,\allowbreak -4,\allowbreak -25,\allowbreak 1,\allowbreak -2)$ & $134.976724\allowbreak{}327872\allowbreak{}125739\allowbreak{}384878\ldots\;\mathrm{MeV}/c^2$\textsuperscript{\(\dagger\)} & $134.976827\allowbreak{}767684\allowbreak{}71\pm 0.000485\allowbreak{}679282\allowbreak{}284233\allowbreak{}51\;\mathrm{MeV}/c^2$\par\footnotesize Catálogo PDG 2026. & Diferencia relativa $-0.766352\allowbreak{}375\,\mathrm{ppm}$; $-0.212979\,\sigma$. Para la firma compuesta declarada, la evaluación se deduce exactamente del carácter beta E3. El selector prospectivo de esa firma es un mapa anterior e independiente de la comparación numérica. \\
% HMT-MASS-EXPLICIT-SIGNATURE:E3B-003
$\pi^\pm$\par firma $(44,\allowbreak 1,\allowbreak -2,\allowbreak 2,\allowbreak -1)$ & $139.570485\allowbreak{}171572\allowbreak{}191374\allowbreak{}734364\ldots\;\mathrm{MeV}/c^2$\textsuperscript{\(\dagger\)} & $139.570390\allowbreak{}983681\allowbreak{}31\pm 0.000182\allowbreak{}007160\allowbreak{}408260\allowbreak{}09\;\mathrm{MeV}/c^2$\par\footnotesize Catálogo PDG 2026. & Diferencia relativa $0.674841\allowbreak{}491\,\mathrm{ppm}$; $0.517495\,\sigma$. Para la firma compuesta declarada, la evaluación se deduce exactamente del carácter beta E3. El selector prospectivo de esa firma es un mapa anterior e independiente de la comparación numérica. \\
% HMT-MASS-EXPLICIT-SIGNATURE:E3B-004
$K^\pm$\par firma $(54,\allowbreak -1,\allowbreak 17,\allowbreak -2,\allowbreak 2)$ & $493.677085\allowbreak{}649530\allowbreak{}612475\allowbreak{}723054\ldots\;\mathrm{MeV}/c^2$\textsuperscript{\(\dagger\)} & $493.676599\allowbreak{}458040\allowbreak{}58\pm 0.015405\allowbreak{}665226\allowbreak{}071509\;\mathrm{MeV}/c^2$\par\footnotesize Catálogo PDG 2026. & Diferencia relativa $0.984838\allowbreak{}030\,\mathrm{ppm}$; $0.031559\,\sigma$. Para la firma compuesta declarada, la evaluación se deduce exactamente del carácter beta E3. El selector prospectivo de esa firma es un mapa anterior e independiente de la comparación numérica. \\
% HMT-MASS-EXPLICIT-SIGNATURE:E3B-005
$K^0$\par firma $(54,\allowbreak 0,\allowbreak 32,\allowbreak 3,\allowbreak -2)$ & $497.611186\allowbreak{}497750\allowbreak{}457358\allowbreak{}558672\ldots\;\mathrm{MeV}/c^2$\textsuperscript{\(\dagger\)} & $497.610767\allowbreak{}531257\allowbreak{}52\pm 0.012990\allowbreak{}199756\allowbreak{}4242\;\mathrm{MeV}/c^2$\par\footnotesize Catálogo PDG 2026. & Diferencia relativa $0.841956\allowbreak{}244\,\mathrm{ppm}$; $0.032252\,\sigma$. Para la firma compuesta declarada, la evaluación se deduce exactamente del carácter beta E3. El selector prospectivo de esa firma es un mapa anterior e independiente de la comparación numérica. \\
% HMT-MASS-EXPLICIT-SIGNATURE:E3B-006
protón $p$\par firma $(59,\allowbreak 0,\allowbreak 9,\allowbreak 1,\allowbreak 0)$ & $938.271751\allowbreak{}546984\allowbreak{}898298\allowbreak{}936787\ldots\;\mathrm{MeV}/c^2$\textsuperscript{\(\dagger\)} & $938.272089\allowbreak{}430000\allowbreak{}05\pm 0.000000\allowbreak{}290000\allowbreak{}000000\allowbreak{}00\;\mathrm{MeV}/c^2$\par\footnotesize Catálogo PDG 2026. & Diferencia relativa $-0.360111\allowbreak{}974\,\mathrm{ppm}$; $-1165.113845\,\sigma$. Para la firma compuesta declarada, la evaluación se deduce exactamente del carácter beta E3. El selector prospectivo de esa firma es un mapa anterior e independiente de la comparación numérica. \\
% HMT-MASS-EXPLICIT-SIGNATURE:E3B-007
neutrón $n$\par firma $(59,\allowbreak 1,\allowbreak -34,\allowbreak 0,\allowbreak 1)$ & $939.565810\allowbreak{}472507\allowbreak{}023312\allowbreak{}664431\ldots\;\mathrm{MeV}/c^2$\textsuperscript{\(\dagger\)} & $939.565421\allowbreak{}939999\allowbreak{}96\pm 0.000000\allowbreak{}480000\allowbreak{}000000\allowbreak{}00\;\mathrm{MeV}/c^2$\par\footnotesize Catálogo PDG 2026. & Diferencia relativa $0.413523\allowbreak{}633\,\mathrm{ppm}$; $809.442723\,\sigma$. Para la firma compuesta declarada, la evaluación se deduce exactamente del carácter beta E3. El selector prospectivo de esa firma es un mapa anterior e independiente de la comparación numérica. \\
\end{longtable}
\endgroup


\Needspace{12\baselineskip}
\paragraph{Invariantes nativos de los sectores topológicos y virtuales.}
Los sectores Fibonacci, Ising, \(\mathbb Z/6\) y las secciones virtuales se comparan
en sus codominios propios. Fibonacci publica el anillo basado y su dimensión de
Frobenius--Perron; los datos \(F/R\) permanecen en la realización posterior. Ising
conserva su anillo de fusión y separa la realización Majorana; \(\mathbb Z/6\)
publica sus caracteres de trenzado; la virtualidad publica horizontes de
persistencia. Esta
tabla mantiene esos invariantes en el mismo mapa de realizaciones sin traducirlos
artificialmente a una unidad de masa.

\begingroup
\setlength{\tabcolsep}{4.5pt}
\renewcommand{\arraystretch}{1.28}
\begin{longtable}{@{}>{\raggedright\arraybackslash}p{.15\textwidth}>{\raggedright\arraybackslash}p{.31\textwidth}>{\raggedright\arraybackslash}p{.25\textwidth}>{\raggedright\arraybackslash}p{.23\textwidth}@{}}
\caption{Sectores topológicos y virtuales: invariantes nativos de fusión, trenzado y persistencia.}\label{tab:masas-invariantes-topologicos-virtuales-rev3}\\
\toprule
\textbf{Sector} & \textbf{Objeto y relaciones} & \textbf{Invariante nativo} & \textbf{Alcance físico} \\
\midrule
\endfirsthead
\toprule
\textbf{Sector} & \textbf{Objeto y relaciones} & \textbf{Invariante nativo} & \textbf{Alcance físico} \\
\midrule
\endhead
\midrule
\multicolumn{4}{r}{\footnotesize Continúa en la página siguiente.}\\
\endfoot
\bottomrule
\endlastfoot
Fibonacci & $K=\mathbb Z1\oplus\mathbb Z\tau$, $\tau^2=1+\tau$. & $\operatorname{FPdim}(\tau)=\varphi$ en el anillo basado. & Las matrices $F/R$, pentágono, hexágono, Hamiltoniano y brecha pertenecen a la realización posterior. \\
Ising & $K=\mathbb Z\{1,\psi,\sigma\}$, $\psi^2=1$, $\psi\sigma=\sigma$, $\sigma^2=1+\psi$. & $\operatorname{FPdim}(\sigma)=\sqrt{2}$; representación de trenzado. & La realización Majorana conserva su propio morfismo físico. \\
$\mathbb Z/6$ & $q_{\epsilon,t}=(-1)^\epsilon\omega_3^t=\zeta_6^k$. & Carácter cuadrático y representación de trenzado sobre los seis sectores. & La observación pertinente es topológica, no una masa universal. \\
Secciones virtuales & Sistema inverso de prolongaciones con horizonte $L(s_N)$. & Registro finito, firma, acarreo y clase de obstrucción. & El observable nativo es la persistencia de la prolongación. \\
\end{longtable}
\endgroup


Los neutrinos permanecen en el codominio de autovalores y mezcla descrito por sus
fichas operatoriales; PMNS realiza el transporte entre las bases de interacción y
masa, mientras la escala pertenece al espectro del operador neutrínico. Las seis
evaluaciones compuestas se obtienen de firmas explícitas mediante el morfismo
general de ligadura hacia la carta escalar. Los sectores topológicos y las secciones virtuales mantienen sus
observables propios de fusión, trenzado y persistencia.

\subsection{Construcción causal y cobertura de las clases de masa}

La arquitectura APP--TRIT--TPK, la ley de masas, el lector bidireccional y
las torres determinan conjuntamente el atlas, los registros cronológicos
enriquecidos, las firmas, el carácter y la clausura basal. Su confluencia en
\eqref{eq:ley-operatorial-general-masas-rectoras} proporciona una única
ecuación de secciones de masa. El dominio completo es el atlas \(81/56/13/324\) y cada realización especifica, según su codominio,
y especifica, según el codominio de cada sector, clasificadores, celdas, rutas,
orientaciones, palabras dodecafásicas, firmas y morfismos de compatibilidad. Los
espectros \(K_D/K_\Omega\) refinan las fibras del atlas, mientras las rutas
cronológicas de \(108\) estados asociadas a \(W\), \(Z\) y \(H\) construyen
directamente sus firmas y caracteres bosónicos.
